%%%%%%%%%%%%%%%%%%%%%%%%%%%%%%%%%%%%%%%%%%%%%%%%%%%%%%%%%%%%
% 14.7 TRIGONOMETRIC IDENTITIES
% The Composition of Advanced Mathematics
%%%%%%%%%%%%%%%%%%%%%%%%%%%%%%%%%%%%%%%%%%%%%%%%%%%%%%%%%%%%

\section{Trigonometric Identities}
\label{sec:trigonometric-identities}

\prerequisite{
Angles, radians, right-triangle trigonometry, the unit circle,
algebraic manipulation, inverse trigonometric functions,
and basic complex-number notation.
}

\learningobjectives{
By the end of this reference section, the reader should be able to:
\begin{itemize}
    \item identify the six fundamental trigonometric functions;
    \item use reciprocal identities;
    \item use quotient identities;
    \item use Pythagorean identities;
    \item recognize even and odd trigonometric functions;
    \item use periodicity identities;
    \item apply cofunction identities;
    \item use reference-angle relationships;
    \item use sum and difference identities;
    \item use double-angle identities;
    \item use half-angle identities;
    \item use power-reduction identities;
    \item use product-to-sum identities;
    \item use sum-to-product identities;
    \item recognize multiple-angle identities;
    \item express trigonometric functions using complex exponentials;
    \item use inverse-trigonometric relationships;
    \item identify principal-value ranges of inverse trigonometric functions;
    \item simplify compositions involving trig and inverse trig functions;
    \item derive selected identities from the unit circle or Euler's formula;
    \item solve identity-verification problems carefully;
    \item recognize common false trigonometric identities;
    \item track domain restrictions when dividing by trigonometric quantities;
    \item connect trigonometric identities with calculus, differential equations,
          Fourier analysis, complex analysis, geometry, and applied mathematics.
\end{itemize}
}

%%%%%%%%%%%%%%%%%%%%%%%%%%%%%%%%%%%%%%%%%%%%%%%%%%%%%%%%%%%%
\subsection{The Six Trigonometric Functions}
\label{subsec:six-trig-functions-reference}
%%%%%%%%%%%%%%%%%%%%%%%%%%%%%%%%%%%%%%%%%%%%%%%%%%%%%%%%%%%%

For an angle \(\theta\), the six basic trigonometric functions are

\[
\boxed{
\sin\theta,
\quad
\cos\theta,
\quad
\tan\theta,
\quad
\csc\theta,
\quad
\sec\theta,
\quad
\cot\theta.
}
\]

On the unit circle,

\[
\boxed{
\cos\theta=x,
\qquad
\sin\theta=y.
}
\]

Thus a point on the unit circle has coordinates

\[
\boxed{
(\cos\theta,\sin\theta).
}
\]

%%%%%%%%%%%%%%%%%%%%%%%%%%%%%%%%%%%%%%%%%%%%%%%%%%%%%%%%%%%%
\subsection{Right-Triangle Definitions}
\label{subsec:right-triangle-trig-definitions}
%%%%%%%%%%%%%%%%%%%%%%%%%%%%%%%%%%%%%%%%%%%%%%%%%%%%%%%%%%%%

For an acute angle \(\theta\) in a right triangle,

\[
\boxed{
\sin\theta
=
\frac{
\text{opposite}
}{
\text{hypotenuse}
},
}
\]

\[
\boxed{
\cos\theta
=
\frac{
\text{adjacent}
}{
\text{hypotenuse}
},
}
\]

and

\[
\boxed{
\tan\theta
=
\frac{
\text{opposite}
}{
\text{adjacent}
}.
}
\]

%%%%%%%%%%%%%%%%%%%%%%%%%%%%%%%%%%%%%%%%%%%%%%%%%%%%%%%%%%%%
\subsection{Reciprocal Identities}
\label{subsec:reciprocal-trig-identities}
%%%%%%%%%%%%%%%%%%%%%%%%%%%%%%%%%%%%%%%%%%%%%%%%%%%%%%%%%%%%

The reciprocal identities are

\[
\boxed{
\csc\theta
=
\frac{1}{\sin\theta},
}
\]

\[
\boxed{
\sec\theta
=
\frac{1}{\cos\theta},
}
\]

and

\[
\boxed{
\cot\theta
=
\frac{1}{\tan\theta}.
}
\]

Equivalently,

\[
\boxed{
\sin\theta
=
\frac{1}{\csc\theta},
}
\]

\[
\boxed{
\cos\theta
=
\frac{1}{\sec\theta},
}
\]

and

\[
\boxed{
\tan\theta
=
\frac{1}{\cot\theta}.
}
\]

%%%%%%%%%%%%%%%%%%%%%%%%%%%%%%%%%%%%%%%%%%%%%%%%%%%%%%%%%%%%
\subsection{Quotient Identities}
\label{subsec:quotient-trig-identities}
%%%%%%%%%%%%%%%%%%%%%%%%%%%%%%%%%%%%%%%%%%%%%%%%%%%%%%%%%%%%

The quotient identities are

\[
\boxed{
\tan\theta
=
\frac{\sin\theta}{\cos\theta},
}
\]

and

\[
\boxed{
\cot\theta
=
\frac{\cos\theta}{\sin\theta}.
}
\]

These are valid wherever the denominators are nonzero.

%%%%%%%%%%%%%%%%%%%%%%%%%%%%%%%%%%%%%%%%%%%%%%%%%%%%%%%%%%%%
\subsection{Primary Pythagorean Identity}
\label{subsec:primary-pythagorean-trig}
%%%%%%%%%%%%%%%%%%%%%%%%%%%%%%%%%%%%%%%%%%%%%%%%%%%%%%%%%%%%

From the unit circle equation

\[
x^2+y^2=1,
\]

with

\[
x=\cos\theta,
\qquad
y=\sin\theta,
\]

we obtain

\[
\boxed{
\sin^2\theta+\cos^2\theta=1.
}
\]

%%%%%%%%%%%%%%%%%%%%%%%%%%%%%%%%%%%%%%%%%%%%%%%%%%%%%%%%%%%%
\subsection{Secondary Pythagorean Identities}
\label{subsec:secondary-pythagorean-trig}
%%%%%%%%%%%%%%%%%%%%%%%%%%%%%%%%%%%%%%%%%%%%%%%%%%%%%%%%%%%%

Divide

\[
\sin^2\theta+\cos^2\theta=1
\]

by \(\cos^2\theta\), where \(\cos\theta\neq0\):

\[
\boxed{
\tan^2\theta+1
=
\sec^2\theta.
}
\]

Divide by \(\sin^2\theta\), where \(\sin\theta\neq0\):

\[
\boxed{
1+\cot^2\theta
=
\csc^2\theta.
}
\]

%%%%%%%%%%%%%%%%%%%%%%%%%%%%%%%%%%%%%%%%%%%%%%%%%%%%%%%%%%%%
\subsection{Pythagorean Rearrangements}
\label{subsec:pythagorean-trig-rearrangements}
%%%%%%%%%%%%%%%%%%%%%%%%%%%%%%%%%%%%%%%%%%%%%%%%%%%%%%%%%%%%

Useful equivalent forms include

\[
\boxed{
1-\sin^2\theta
=
\cos^2\theta,
}
\]

\[
\boxed{
1-\cos^2\theta
=
\sin^2\theta,
}
\]

\[
\boxed{
\sec^2\theta-\tan^2\theta
=
1,
}
\]

and

\[
\boxed{
\csc^2\theta-\cot^2\theta
=
1.
}
\]

%%%%%%%%%%%%%%%%%%%%%%%%%%%%%%%%%%%%%%%%%%%%%%%%%%%%%%%%%%%%
\subsection{Even and Odd Identities}
\label{subsec:even-odd-trig-identities}
%%%%%%%%%%%%%%%%%%%%%%%%%%%%%%%%%%%%%%%%%%%%%%%%%%%%%%%%%%%%

Cosine and secant are even:

\[
\boxed{
\cos(-\theta)
=
\cos\theta,
}
\]

\[
\boxed{
\sec(-\theta)
=
\sec\theta.
}
\]

Sine, tangent, cosecant, and cotangent are odd:

\[
\boxed{
\sin(-\theta)
=
-\sin\theta,
}
\]

\[
\boxed{
\tan(-\theta)
=
-\tan\theta,
}
\]

\[
\boxed{
\csc(-\theta)
=
-\csc\theta,
}
\]

\[
\boxed{
\cot(-\theta)
=
-\cot\theta.
}
\]

%%%%%%%%%%%%%%%%%%%%%%%%%%%%%%%%%%%%%%%%%%%%%%%%%%%%%%%%%%%%
\subsection{Periodicity}
\label{subsec:trig-periodicity-reference}
%%%%%%%%%%%%%%%%%%%%%%%%%%%%%%%%%%%%%%%%%%%%%%%%%%%%%%%%%%%%

Sine and cosine have period \(2\pi\):

\[
\boxed{
\sin(\theta+2\pi)
=
\sin\theta,
}
\]

\[
\boxed{
\cos(\theta+2\pi)
=
\cos\theta.
}
\]

Therefore,

\[
\boxed{
\sec(\theta+2\pi)
=
\sec\theta,
}
\]

and

\[
\boxed{
\csc(\theta+2\pi)
=
\csc\theta.
}
\]

%%%%%%%%%%%%%%%%%%%%%%%%%%%%%%%%%%%%%%%%%%%%%%%%%%%%%%%%%%%%
\subsection{Tangent and Cotangent Periodicity}
\label{subsec:tangent-cotangent-period}
%%%%%%%%%%%%%%%%%%%%%%%%%%%%%%%%%%%%%%%%%%%%%%%%%%%%%%%%%%%%

Tangent and cotangent have period \(\pi\):

\[
\boxed{
\tan(\theta+\pi)
=
\tan\theta,
}
\]

and

\[
\boxed{
\cot(\theta+\pi)
=
\cot\theta.
}
\]

%%%%%%%%%%%%%%%%%%%%%%%%%%%%%%%%%%%%%%%%%%%%%%%%%%%%%%%%%%%%
\subsection{Cofunction Identities}
\label{subsec:cofunction-identities-reference}
%%%%%%%%%%%%%%%%%%%%%%%%%%%%%%%%%%%%%%%%%%%%%%%%%%%%%%%%%%%%

For complementary angles,

\[
\boxed{
\sin\left(
\frac{\pi}{2}-\theta
\right)
=
\cos\theta,
}
\]

\[
\boxed{
\cos\left(
\frac{\pi}{2}-\theta
\right)
=
\sin\theta,
}
\]

\[
\boxed{
\tan\left(
\frac{\pi}{2}-\theta
\right)
=
\cot\theta,
}
\]

\[
\boxed{
\cot\left(
\frac{\pi}{2}-\theta
\right)
=
\tan\theta,
}
\]

\[
\boxed{
\sec\left(
\frac{\pi}{2}-\theta
\right)
=
\csc\theta,
}
\]

and

\[
\boxed{
\csc\left(
\frac{\pi}{2}-\theta
\right)
=
\sec\theta.
}
\]

%%%%%%%%%%%%%%%%%%%%%%%%%%%%%%%%%%%%%%%%%%%%%%%%%%%%%%%%%%%%
\subsection{Reference-Angle Sign Patterns}
\label{subsec:reference-angle-sign-patterns}
%%%%%%%%%%%%%%%%%%%%%%%%%%%%%%%%%%%%%%%%%%%%%%%%%%%%%%%%%%%%

The signs of sine and cosine determine the signs of all six trigonometric
functions.

A common quadrant mnemonic is:

\[
\boxed{
\begin{array}{c|c}
\text{Quadrant} & \text{Positive Functions} \\ \hline
I & \text{all} \\
II & \sin,\csc \\
III & \tan,\cot \\
IV & \cos,\sec
\end{array}
}
\]

%%%%%%%%%%%%%%%%%%%%%%%%%%%%%%%%%%%%%%%%%%%%%%%%%%%%%%%%%%%%
\subsection{Quadrant Reduction Identities}
\label{subsec:quadrant-reduction-identities}
%%%%%%%%%%%%%%%%%%%%%%%%%%%%%%%%%%%%%%%%%%%%%%%%%%%%%%%%%%%%

Useful identities include

\[
\boxed{
\sin(\pi-\theta)
=
\sin\theta,
}
\]

\[
\boxed{
\cos(\pi-\theta)
=
-\cos\theta,
}
\]

\[
\boxed{
\tan(\pi-\theta)
=
-\tan\theta,
}
\]

\[
\boxed{
\sin(\pi+\theta)
=
-\sin\theta,
}
\]

\[
\boxed{
\cos(\pi+\theta)
=
-\cos\theta,
}
\]

and

\[
\boxed{
\tan(\pi+\theta)
=
\tan\theta.
}
\]

%%%%%%%%%%%%%%%%%%%%%%%%%%%%%%%%%%%%%%%%%%%%%%%%%%%%%%%%%%%%
\subsection{Additional Reduction Identities}
\label{subsec:additional-reduction-identities}
%%%%%%%%%%%%%%%%%%%%%%%%%%%%%%%%%%%%%%%%%%%%%%%%%%%%%%%%%%%%

Also,

\[
\boxed{
\sin(2\pi-\theta)
=
-\sin\theta,
}
\]

\[
\boxed{
\cos(2\pi-\theta)
=
\cos\theta,
}
\]

and

\[
\boxed{
\tan(2\pi-\theta)
=
-\tan\theta.
}
\]

%%%%%%%%%%%%%%%%%%%%%%%%%%%%%%%%%%%%%%%%%%%%%%%%%%%%%%%%%%%%
\subsection{Sine Sum Identity}
\label{subsec:sine-sum-identity}
%%%%%%%%%%%%%%%%%%%%%%%%%%%%%%%%%%%%%%%%%%%%%%%%%%%%%%%%%%%%

\[
\boxed{
\sin(\alpha+\beta)
=
\sin\alpha\cos\beta
+
\cos\alpha\sin\beta.
}
\]

%%%%%%%%%%%%%%%%%%%%%%%%%%%%%%%%%%%%%%%%%%%%%%%%%%%%%%%%%%%%
\subsection{Sine Difference Identity}
\label{subsec:sine-difference-identity}
%%%%%%%%%%%%%%%%%%%%%%%%%%%%%%%%%%%%%%%%%%%%%%%%%%%%%%%%%%%%

\[
\boxed{
\sin(\alpha-\beta)
=
\sin\alpha\cos\beta
-
\cos\alpha\sin\beta.
}
\]

%%%%%%%%%%%%%%%%%%%%%%%%%%%%%%%%%%%%%%%%%%%%%%%%%%%%%%%%%%%%
\subsection{Cosine Sum Identity}
\label{subsec:cosine-sum-identity}
%%%%%%%%%%%%%%%%%%%%%%%%%%%%%%%%%%%%%%%%%%%%%%%%%%%%%%%%%%%%

\[
\boxed{
\cos(\alpha+\beta)
=
\cos\alpha\cos\beta
-
\sin\alpha\sin\beta.
}
\]

%%%%%%%%%%%%%%%%%%%%%%%%%%%%%%%%%%%%%%%%%%%%%%%%%%%%%%%%%%%%
\subsection{Cosine Difference Identity}
\label{subsec:cosine-difference-identity}
%%%%%%%%%%%%%%%%%%%%%%%%%%%%%%%%%%%%%%%%%%%%%%%%%%%%%%%%%%%%

\[
\boxed{
\cos(\alpha-\beta)
=
\cos\alpha\cos\beta
+
\sin\alpha\sin\beta.
}
\]

%%%%%%%%%%%%%%%%%%%%%%%%%%%%%%%%%%%%%%%%%%%%%%%%%%%%%%%%%%%%
\subsection{Tangent Sum Identity}
\label{subsec:tangent-sum-identity}
%%%%%%%%%%%%%%%%%%%%%%%%%%%%%%%%%%%%%%%%%%%%%%%%%%%%%%%%%%%%

Where defined,

\[
\boxed{
\tan(\alpha+\beta)
=
\frac{
\tan\alpha+\tan\beta
}{
1-\tan\alpha\tan\beta
}.
}
\]

%%%%%%%%%%%%%%%%%%%%%%%%%%%%%%%%%%%%%%%%%%%%%%%%%%%%%%%%%%%%
\subsection{Tangent Difference Identity}
\label{subsec:tangent-difference-identity}
%%%%%%%%%%%%%%%%%%%%%%%%%%%%%%%%%%%%%%%%%%%%%%%%%%%%%%%%%%%%

Where defined,

\[
\boxed{
\tan(\alpha-\beta)
=
\frac{
\tan\alpha-\tan\beta
}{
1+\tan\alpha\tan\beta
}.
}
\]

%%%%%%%%%%%%%%%%%%%%%%%%%%%%%%%%%%%%%%%%%%%%%%%%%%%%%%%%%%%%
\subsection{Worked Example: Exact Sine Value}
\label{subsec:worked-sine-sum}
%%%%%%%%%%%%%%%%%%%%%%%%%%%%%%%%%%%%%%%%%%%%%%%%%%%%%%%%%%%%

\begin{workedexamplebox}{Find \(\sin75^\circ\)}

Write

\[
75^\circ
=
45^\circ+30^\circ.
\]

Then

\[
\sin75^\circ
=
\sin45^\circ\cos30^\circ
+
\cos45^\circ\sin30^\circ.
\]

Using

\[
\sin45^\circ
=
\cos45^\circ
=
\frac{\sqrt2}{2},
\]

\[
\cos30^\circ
=
\frac{\sqrt3}{2},
\]

and

\[
\sin30^\circ
=
\frac12,
\]

we obtain

\[
\sin75^\circ
=
\frac{\sqrt6}{4}
+
\frac{\sqrt2}{4}.
\]

\begin{answerbox}
\[
\boxed{
\sin75^\circ
=
\frac{
\sqrt6+\sqrt2
}{4}.
}
\]
\end{answerbox}

\end{workedexamplebox}

%%%%%%%%%%%%%%%%%%%%%%%%%%%%%%%%%%%%%%%%%%%%%%%%%%%%%%%%%%%%
\subsection{Worked Example: Exact Cosine Value}
\label{subsec:worked-cosine-difference}
%%%%%%%%%%%%%%%%%%%%%%%%%%%%%%%%%%%%%%%%%%%%%%%%%%%%%%%%%%%%

\begin{workedexamplebox}{Find \(\cos15^\circ\)}

Write

\[
15^\circ
=
45^\circ-30^\circ.
\]

Then

\[
\cos15^\circ
=
\cos45^\circ\cos30^\circ
+
\sin45^\circ\sin30^\circ.
\]

Therefore,

\begin{answerbox}
\[
\boxed{
\cos15^\circ
=
\frac{
\sqrt6+\sqrt2
}{4}.
}
\]
\end{answerbox}

\end{workedexamplebox}

%%%%%%%%%%%%%%%%%%%%%%%%%%%%%%%%%%%%%%%%%%%%%%%%%%%%%%%%%%%%
\subsection{Double-Angle Identity for Sine}
\label{subsec:sine-double-angle}
%%%%%%%%%%%%%%%%%%%%%%%%%%%%%%%%%%%%%%%%%%%%%%%%%%%%%%%%%%%%

Set

\[
\alpha=\beta=\theta
\]

in the sine sum identity:

\[
\boxed{
\sin2\theta
=
2\sin\theta\cos\theta.
}
\]

%%%%%%%%%%%%%%%%%%%%%%%%%%%%%%%%%%%%%%%%%%%%%%%%%%%%%%%%%%%%
\subsection{Double-Angle Identity for Cosine}
\label{subsec:cosine-double-angle}
%%%%%%%%%%%%%%%%%%%%%%%%%%%%%%%%%%%%%%%%%%%%%%%%%%%%%%%%%%%%

From the cosine sum identity,

\[
\boxed{
\cos2\theta
=
\cos^2\theta-\sin^2\theta.
}
\]

Using

\[
\sin^2\theta+\cos^2\theta=1,
\]

we also obtain

\[
\boxed{
\cos2\theta
=
2\cos^2\theta-1,
}
\]

and

\[
\boxed{
\cos2\theta
=
1-2\sin^2\theta.
}
\]

%%%%%%%%%%%%%%%%%%%%%%%%%%%%%%%%%%%%%%%%%%%%%%%%%%%%%%%%%%%%
\subsection{Double-Angle Identity for Tangent}
\label{subsec:tangent-double-angle}
%%%%%%%%%%%%%%%%%%%%%%%%%%%%%%%%%%%%%%%%%%%%%%%%%%%%%%%%%%%%

Where defined,

\[
\boxed{
\tan2\theta
=
\frac{
2\tan\theta
}{
1-\tan^2\theta
}.
}
\]

%%%%%%%%%%%%%%%%%%%%%%%%%%%%%%%%%%%%%%%%%%%%%%%%%%%%%%%%%%%%
\subsection{Power-Reduction Identity for Sine}
\label{subsec:sine-power-reduction}
%%%%%%%%%%%%%%%%%%%%%%%%%%%%%%%%%%%%%%%%%%%%%%%%%%%%%%%%%%%%

From

\[
\cos2\theta
=
1-2\sin^2\theta,
\]

we obtain

\[
\boxed{
\sin^2\theta
=
\frac{
1-\cos2\theta
}{2}.
}
\]

%%%%%%%%%%%%%%%%%%%%%%%%%%%%%%%%%%%%%%%%%%%%%%%%%%%%%%%%%%%%
\subsection{Power-Reduction Identity for Cosine}
\label{subsec:cosine-power-reduction}
%%%%%%%%%%%%%%%%%%%%%%%%%%%%%%%%%%%%%%%%%%%%%%%%%%%%%%%%%%%%

From

\[
\cos2\theta
=
2\cos^2\theta-1,
\]

we obtain

\[
\boxed{
\cos^2\theta
=
\frac{
1+\cos2\theta
}{2}.
}
\]

%%%%%%%%%%%%%%%%%%%%%%%%%%%%%%%%%%%%%%%%%%%%%%%%%%%%%%%%%%%%
\subsection{Half-Angle Identity for Sine}
\label{subsec:sine-half-angle}
%%%%%%%%%%%%%%%%%%%%%%%%%%%%%%%%%%%%%%%%%%%%%%%%%%%%%%%%%%%%

Replacing \(\theta\) by \(\theta/2\),

\[
\boxed{
\sin^2
\left(
\frac{\theta}{2}
\right)
=
\frac{
1-\cos\theta
}{2}.
}
\]

Therefore,

\[
\boxed{
\sin
\left(
\frac{\theta}{2}
\right)
=
\pm
\sqrt{
\frac{
1-\cos\theta
}{2}
}.
}
\]

The sign depends on the quadrant of \(\theta/2\).

%%%%%%%%%%%%%%%%%%%%%%%%%%%%%%%%%%%%%%%%%%%%%%%%%%%%%%%%%%%%
\subsection{Half-Angle Identity for Cosine}
\label{subsec:cosine-half-angle}
%%%%%%%%%%%%%%%%%%%%%%%%%%%%%%%%%%%%%%%%%%%%%%%%%%%%%%%%%%%%

\[
\boxed{
\cos^2
\left(
\frac{\theta}{2}
\right)
=
\frac{
1+\cos\theta
}{2}.
}
\]

Thus

\[
\boxed{
\cos
\left(
\frac{\theta}{2}
\right)
=
\pm
\sqrt{
\frac{
1+\cos\theta
}{2}
}.
}
\]

%%%%%%%%%%%%%%%%%%%%%%%%%%%%%%%%%%%%%%%%%%%%%%%%%%%%%%%%%%%%
\subsection{Half-Angle Identity for Tangent}
\label{subsec:tangent-half-angle}
%%%%%%%%%%%%%%%%%%%%%%%%%%%%%%%%%%%%%%%%%%%%%%%%%%%%%%%%%%%%

Useful forms include

\[
\boxed{
\tan
\left(
\frac{\theta}{2}
\right)
=
\frac{
\sin\theta
}{
1+\cos\theta
},
}
\]

and

\[
\boxed{
\tan
\left(
\frac{\theta}{2}
\right)
=
\frac{
1-\cos\theta
}{
\sin\theta
}.
}
\]

Where square-root form is appropriate,

\[
\boxed{
\tan
\left(
\frac{\theta}{2}
\right)
=
\pm
\sqrt{
\frac{
1-\cos\theta
}{
1+\cos\theta
}
}.
}
\]

%%%%%%%%%%%%%%%%%%%%%%%%%%%%%%%%%%%%%%%%%%%%%%%%%%%%%%%%%%%%
\subsection{Triple-Angle Identity for Sine}
\label{subsec:sine-triple-angle}
%%%%%%%%%%%%%%%%%%%%%%%%%%%%%%%%%%%%%%%%%%%%%%%%%%%%%%%%%%%%

\[
\boxed{
\sin3\theta
=
3\sin\theta
-
4\sin^3\theta.
}
\]

%%%%%%%%%%%%%%%%%%%%%%%%%%%%%%%%%%%%%%%%%%%%%%%%%%%%%%%%%%%%
\subsection{Triple-Angle Identity for Cosine}
\label{subsec:cosine-triple-angle}
%%%%%%%%%%%%%%%%%%%%%%%%%%%%%%%%%%%%%%%%%%%%%%%%%%%%%%%%%%%%

\[
\boxed{
\cos3\theta
=
4\cos^3\theta
-
3\cos\theta.
}
\]

%%%%%%%%%%%%%%%%%%%%%%%%%%%%%%%%%%%%%%%%%%%%%%%%%%%%%%%%%%%%
\subsection{Triple-Angle Identity for Tangent}
\label{subsec:tangent-triple-angle}
%%%%%%%%%%%%%%%%%%%%%%%%%%%%%%%%%%%%%%%%%%%%%%%%%%%%%%%%%%%%

Where defined,

\[
\boxed{
\tan3\theta
=
\frac{
3\tan\theta-\tan^3\theta
}{
1-3\tan^2\theta
}.
}
\]

%%%%%%%%%%%%%%%%%%%%%%%%%%%%%%%%%%%%%%%%%%%%%%%%%%%%%%%%%%%%
\subsection{Product-to-Sum: Sine-Sine}
\label{subsec:product-to-sum-sine-sine}
%%%%%%%%%%%%%%%%%%%%%%%%%%%%%%%%%%%%%%%%%%%%%%%%%%%%%%%%%%%%

\[
\boxed{
\sin\alpha\sin\beta
=
\frac12
\left[
\cos(\alpha-\beta)
-
\cos(\alpha+\beta)
\right].
}
\]

%%%%%%%%%%%%%%%%%%%%%%%%%%%%%%%%%%%%%%%%%%%%%%%%%%%%%%%%%%%%
\subsection{Product-to-Sum: Cosine-Cosine}
\label{subsec:product-to-sum-cos-cos}
%%%%%%%%%%%%%%%%%%%%%%%%%%%%%%%%%%%%%%%%%%%%%%%%%%%%%%%%%%%%

\[
\boxed{
\cos\alpha\cos\beta
=
\frac12
\left[
\cos(\alpha-\beta)
+
\cos(\alpha+\beta)
\right].
}
\]

%%%%%%%%%%%%%%%%%%%%%%%%%%%%%%%%%%%%%%%%%%%%%%%%%%%%%%%%%%%%
\subsection{Product-to-Sum: Sine-Cosine}
\label{subsec:product-to-sum-sine-cosine}
%%%%%%%%%%%%%%%%%%%%%%%%%%%%%%%%%%%%%%%%%%%%%%%%%%%%%%%%%%%%

\[
\boxed{
\sin\alpha\cos\beta
=
\frac12
\left[
\sin(\alpha+\beta)
+
\sin(\alpha-\beta)
\right].
}
\]

%%%%%%%%%%%%%%%%%%%%%%%%%%%%%%%%%%%%%%%%%%%%%%%%%%%%%%%%%%%%
\subsection{Product-to-Sum: Cosine-Sine}
\label{subsec:product-to-sum-cosine-sine}
%%%%%%%%%%%%%%%%%%%%%%%%%%%%%%%%%%%%%%%%%%%%%%%%%%%%%%%%%%%%

\[
\boxed{
\cos\alpha\sin\beta
=
\frac12
\left[
\sin(\alpha+\beta)
-
\sin(\alpha-\beta)
\right].
}
\]

%%%%%%%%%%%%%%%%%%%%%%%%%%%%%%%%%%%%%%%%%%%%%%%%%%%%%%%%%%%%
\subsection{Sum-to-Product: Sine Sum}
\label{subsec:sum-to-product-sine-sum}
%%%%%%%%%%%%%%%%%%%%%%%%%%%%%%%%%%%%%%%%%%%%%%%%%%%%%%%%%%%%

\[
\boxed{
\sin\alpha+\sin\beta
=
2
\sin
\left(
\frac{\alpha+\beta}{2}
\right)
\cos
\left(
\frac{\alpha-\beta}{2}
\right).
}
\]

%%%%%%%%%%%%%%%%%%%%%%%%%%%%%%%%%%%%%%%%%%%%%%%%%%%%%%%%%%%%
\subsection{Sum-to-Product: Sine Difference}
\label{subsec:sum-to-product-sine-difference}
%%%%%%%%%%%%%%%%%%%%%%%%%%%%%%%%%%%%%%%%%%%%%%%%%%%%%%%%%%%%

\[
\boxed{
\sin\alpha-\sin\beta
=
2
\cos
\left(
\frac{\alpha+\beta}{2}
\right)
\sin
\left(
\frac{\alpha-\beta}{2}
\right).
}
\]

%%%%%%%%%%%%%%%%%%%%%%%%%%%%%%%%%%%%%%%%%%%%%%%%%%%%%%%%%%%%
\subsection{Sum-to-Product: Cosine Sum}
\label{subsec:sum-to-product-cosine-sum}
%%%%%%%%%%%%%%%%%%%%%%%%%%%%%%%%%%%%%%%%%%%%%%%%%%%%%%%%%%%%

\[
\boxed{
\cos\alpha+\cos\beta
=
2
\cos
\left(
\frac{\alpha+\beta}{2}
\right)
\cos
\left(
\frac{\alpha-\beta}{2}
\right).
}
\]

%%%%%%%%%%%%%%%%%%%%%%%%%%%%%%%%%%%%%%%%%%%%%%%%%%%%%%%%%%%%
\subsection{Sum-to-Product: Cosine Difference}
\label{subsec:sum-to-product-cosine-difference}
%%%%%%%%%%%%%%%%%%%%%%%%%%%%%%%%%%%%%%%%%%%%%%%%%%%%%%%%%%%%

\[
\boxed{
\cos\alpha-\cos\beta
=
-2
\sin
\left(
\frac{\alpha+\beta}{2}
\right)
\sin
\left(
\frac{\alpha-\beta}{2}
\right).
}
\]

%%%%%%%%%%%%%%%%%%%%%%%%%%%%%%%%%%%%%%%%%%%%%%%%%%%%%%%%%%%%
\subsection{Worked Example: Product-to-Sum}
\label{subsec:worked-product-to-sum}
%%%%%%%%%%%%%%%%%%%%%%%%%%%%%%%%%%%%%%%%%%%%%%%%%%%%%%%%%%%%

\begin{workedexamplebox}{Rewrite a Product}

Rewrite

\[
\sin3x\cos2x
\]

as a sum.

Use

\[
\sin\alpha\cos\beta
=
\frac12
\left[
\sin(\alpha+\beta)
+
\sin(\alpha-\beta)
\right].
\]

Thus

\[
\sin3x\cos2x
=
\frac12
\left[
\sin5x+\sin x
\right].
\]

\begin{answerbox}
\[
\boxed{
\sin3x\cos2x
=
\frac12
(\sin5x+\sin x).
}
\]
\end{answerbox}

\end{workedexamplebox}

%%%%%%%%%%%%%%%%%%%%%%%%%%%%%%%%%%%%%%%%%%%%%%%%%%%%%%%%%%%%
\subsection{Worked Example: Sum-to-Product}
\label{subsec:worked-sum-to-product}
%%%%%%%%%%%%%%%%%%%%%%%%%%%%%%%%%%%%%%%%%%%%%%%%%%%%%%%%%%%%

\begin{workedexamplebox}{Rewrite a Sum}

Rewrite

\[
\cos5x+\cos x.
\]

Using

\[
\cos\alpha+\cos\beta
=
2
\cos
\left(
\frac{\alpha+\beta}{2}
\right)
\cos
\left(
\frac{\alpha-\beta}{2}
\right),
\]

we obtain

\[
\cos5x+\cos x
=
2\cos3x\cos2x.
\]

\begin{answerbox}
\[
\boxed{
\cos5x+\cos x
=
2\cos3x\cos2x.
}
\]
\end{answerbox}

\end{workedexamplebox}

%%%%%%%%%%%%%%%%%%%%%%%%%%%%%%%%%%%%%%%%%%%%%%%%%%%%%%%%%%%%
\subsection{Tangent in Terms of Sine and Cosine}
\label{subsec:tangent-sine-cos-reference}
%%%%%%%%%%%%%%%%%%%%%%%%%%%%%%%%%%%%%%%%%%%%%%%%%%%%%%%%%%%%

A useful simplification strategy is

\[
\boxed{
\tan\theta
=
\frac{\sin\theta}{\cos\theta}.
}
\]

Similarly,

\[
\boxed{
\cot\theta
=
\frac{\cos\theta}{\sin\theta}.
}
\]

Many identity-verification problems become easier after converting all
functions to sine and cosine.

%%%%%%%%%%%%%%%%%%%%%%%%%%%%%%%%%%%%%%%%%%%%%%%%%%%%%%%%%%%%
\subsection{Secant and Cosecant in Terms of Sine and Cosine}
\label{subsec:sec-csc-sine-cos}
%%%%%%%%%%%%%%%%%%%%%%%%%%%%%%%%%%%%%%%%%%%%%%%%%%%%%%%%%%%%

\[
\boxed{
\sec\theta
=
\frac{1}{\cos\theta},
}
\]

and

\[
\boxed{
\csc\theta
=
\frac{1}{\sin\theta}.
}
\]

Thus every elementary trigonometric expression can be rewritten using only
sine and cosine.

%%%%%%%%%%%%%%%%%%%%%%%%%%%%%%%%%%%%%%%%%%%%%%%%%%%%%%%%%%%%
\subsection{Useful Identity With Tangent and Secant}
\label{subsec:tangent-secant-factor-identity}
%%%%%%%%%%%%%%%%%%%%%%%%%%%%%%%%%%%%%%%%%%%%%%%%%%%%%%%%%%%%

Since

\[
\sec^2\theta-\tan^2\theta=1,
\]

we can factor:

\[
\boxed{
(\sec\theta-\tan\theta)
(\sec\theta+\tan\theta)
=
1.
}
\]

Therefore, where defined,

\[
\boxed{
\sec\theta-\tan\theta
=
\frac{
1
}{
\sec\theta+\tan\theta
}.
}
\]

%%%%%%%%%%%%%%%%%%%%%%%%%%%%%%%%%%%%%%%%%%%%%%%%%%%%%%%%%%%%
\subsection{Useful Identity With Cosecant and Cotangent}
\label{subsec:csc-cot-factor-identity}
%%%%%%%%%%%%%%%%%%%%%%%%%%%%%%%%%%%%%%%%%%%%%%%%%%%%%%%%%%%%

Similarly,

\[
\boxed{
(\csc\theta-\cot\theta)
(\csc\theta+\cot\theta)
=
1.
}
\]

Hence,

\[
\boxed{
\csc\theta-\cot\theta
=
\frac{
1
}{
\csc\theta+\cot\theta
}
}
\]

where defined.

%%%%%%%%%%%%%%%%%%%%%%%%%%%%%%%%%%%%%%%%%%%%%%%%%%%%%%%%%%%%
\subsection{Auxiliary-Angle Form}
\label{subsec:auxiliary-angle-reference}
%%%%%%%%%%%%%%%%%%%%%%%%%%%%%%%%%%%%%%%%%%%%%%%%%%%%%%%%%%%%

An expression

\[
a\cos x+b\sin x
\]

can be rewritten as

\[
\boxed{
R\cos(x-\phi),
}
\]

where

\[
\boxed{
R
=
\sqrt{a^2+b^2}.
}
\]

Choose \(\phi\) such that

\[
\boxed{
\cos\phi
=
\frac{a}{R},
\qquad
\sin\phi
=
\frac{b}{R}.
}
\]

%%%%%%%%%%%%%%%%%%%%%%%%%%%%%%%%%%%%%%%%%%%%%%%%%%%%%%%%%%%%
\subsection{Alternative Auxiliary-Angle Form}
\label{subsec:auxiliary-angle-sine-form}
%%%%%%%%%%%%%%%%%%%%%%%%%%%%%%%%%%%%%%%%%%%%%%%%%%%%%%%%%%%%

Likewise,

\[
a\sin x+b\cos x
\]

can be written

\[
\boxed{
R\sin(x+\phi),
}
\]

with

\[
R=\sqrt{a^2+b^2}
\]

and suitable \(\phi\).

%%%%%%%%%%%%%%%%%%%%%%%%%%%%%%%%%%%%%%%%%%%%%%%%%%%%%%%%%%%%
\subsection{Worked Example: Combining Sine and Cosine}
\label{subsec:worked-combining-sine-cosine}
%%%%%%%%%%%%%%%%%%%%%%%%%%%%%%%%%%%%%%%%%%%%%%%%%%%%%%%%%%%%

\begin{workedexamplebox}{Rewrite \(3\cos x+4\sin x\)}

Let

\[
R
=
\sqrt{3^2+4^2}
=
5.
\]

Choose \(\phi\) such that

\[
\cos\phi
=
\frac35,
\qquad
\sin\phi
=
\frac45.
\]

Then

\[
5\cos(x-\phi)
=
5
\left(
\cos x\cos\phi
+
\sin x\sin\phi
\right).
\]

Thus

\[
5\cos(x-\phi)
=
3\cos x+4\sin x.
\]

\begin{answerbox}
\[
\boxed{
3\cos x+4\sin x
=
5\cos(x-\phi),
}
\]

where

\[
\boxed{
\phi
=
\arctan\left(\frac43\right).
}
\]
\end{answerbox}

\end{workedexamplebox}

%%%%%%%%%%%%%%%%%%%%%%%%%%%%%%%%%%%%%%%%%%%%%%%%%%%%%%%%%%%%
\subsection{Maximum of a Linear Sine-Cosine Combination}
\label{subsec:max-linear-trig-combination}
%%%%%%%%%%%%%%%%%%%%%%%%%%%%%%%%%%%%%%%%%%%%%%%%%%%%%%%%%%%%

Since

\[
a\cos x+b\sin x
=
R\cos(x-\phi),
\]

where

\[
R=\sqrt{a^2+b^2},
\]

we have

\[
\boxed{
-\sqrt{a^2+b^2}
\leq
a\cos x+b\sin x
\leq
\sqrt{a^2+b^2}.
}
\]

Therefore,

\[
\boxed{
\max_x
(a\cos x+b\sin x)
=
\sqrt{a^2+b^2}.
}
\]

%%%%%%%%%%%%%%%%%%%%%%%%%%%%%%%%%%%%%%%%%%%%%%%%%%%%%%%%%%%%
\subsection{Complex Exponential Identities}
\label{subsec:complex-exponential-trig}
%%%%%%%%%%%%%%%%%%%%%%%%%%%%%%%%%%%%%%%%%%%%%%%%%%%%%%%%%%%%

Euler's formula is

\[
\boxed{
e^{i\theta}
=
\cos\theta+i\sin\theta.
}
\]

Its conjugate is

\[
\boxed{
e^{-i\theta}
=
\cos\theta-i\sin\theta.
}
\]

Adding gives

\[
\boxed{
\cos\theta
=
\frac{
e^{i\theta}
+
e^{-i\theta}
}{2}.
}
\]

Subtracting gives

\[
\boxed{
\sin\theta
=
\frac{
e^{i\theta}
-
e^{-i\theta}
}{
2i
}.
}
\]

%%%%%%%%%%%%%%%%%%%%%%%%%%%%%%%%%%%%%%%%%%%%%%%%%%%%%%%%%%%%
\subsection{De Moivre's Formula}
\label{subsec:de-moivre-trig-reference}
%%%%%%%%%%%%%%%%%%%%%%%%%%%%%%%%%%%%%%%%%%%%%%%%%%%%%%%%%%%%

For integer \(n\),

\[
\boxed{
(\cos\theta+i\sin\theta)^n
=
\cos(n\theta)
+
i\sin(n\theta).
}
\]

Equivalently,

\[
\boxed{
e^{in\theta}
=
\cos(n\theta)
+
i\sin(n\theta).
}
\]

This provides a systematic way to derive multiple-angle identities.

%%%%%%%%%%%%%%%%%%%%%%%%%%%%%%%%%%%%%%%%%%%%%%%%%%%%%%%%%%%%
\subsection{Chebyshev Connection}
\label{subsec:chebyshev-trig-connection}
%%%%%%%%%%%%%%%%%%%%%%%%%%%%%%%%%%%%%%%%%%%%%%%%%%%%%%%%%%%%

For the Chebyshev polynomial \(T_n\),

\[
\boxed{
T_n(\cos\theta)
=
\cos(n\theta).
}
\]

Thus multiple-angle cosine formulas can be expressed polynomially in
\(\cos\theta\).

For example,

\[
\boxed{
T_3(x)
=
4x^3-3x.
}
\]

%%%%%%%%%%%%%%%%%%%%%%%%%%%%%%%%%%%%%%%%%%%%%%%%%%%%%%%%%%%%
\subsection{Sine Multiple-Angle Connection}
\label{subsec:sine-chebyshev-connection}
%%%%%%%%%%%%%%%%%%%%%%%%%%%%%%%%%%%%%%%%%%%%%%%%%%%%%%%%%%%%

Using Chebyshev polynomials of the second kind,

\[
\boxed{
\sin((n+1)\theta)
=
\sin\theta\,
U_n(\cos\theta).
}
\]

This connects trigonometric identities with orthogonal polynomials.

%%%%%%%%%%%%%%%%%%%%%%%%%%%%%%%%%%%%%%%%%%%%%%%%%%%%%%%%%%%%
\subsection{Inverse Sine}
\label{subsec:inverse-sine-reference}
%%%%%%%%%%%%%%%%%%%%%%%%%%%%%%%%%%%%%%%%%%%%%%%%%%%%%%%%%%%%

The inverse sine function is

\[
\boxed{
y=\arcsin x
}
\]

with principal range

\[
\boxed{
-\frac{\pi}{2}
\leq
y
\leq
\frac{\pi}{2}.
}
\]

Its domain is

\[
\boxed{
-1\leq x\leq1.
}
\]

%%%%%%%%%%%%%%%%%%%%%%%%%%%%%%%%%%%%%%%%%%%%%%%%%%%%%%%%%%%%
\subsection{Inverse Cosine}
\label{subsec:inverse-cosine-reference}
%%%%%%%%%%%%%%%%%%%%%%%%%%%%%%%%%%%%%%%%%%%%%%%%%%%%%%%%%%%%

The inverse cosine function is

\[
\boxed{
y=\arccos x
}
\]

with principal range

\[
\boxed{
0\leq y\leq\pi.
}
\]

Its domain is

\[
\boxed{
-1\leq x\leq1.
}
\]

%%%%%%%%%%%%%%%%%%%%%%%%%%%%%%%%%%%%%%%%%%%%%%%%%%%%%%%%%%%%
\subsection{Inverse Tangent}
\label{subsec:inverse-tangent-reference}
%%%%%%%%%%%%%%%%%%%%%%%%%%%%%%%%%%%%%%%%%%%%%%%%%%%%%%%%%%%%

The inverse tangent function is

\[
\boxed{
y=\arctan x
}
\]

with principal range

\[
\boxed{
-\frac{\pi}{2}
<
y
<
\frac{\pi}{2}.
}
\]

Its domain is

\[
\boxed{
x\in\mathbb{R}.
}
\]

%%%%%%%%%%%%%%%%%%%%%%%%%%%%%%%%%%%%%%%%%%%%%%%%%%%%%%%%%%%%
\subsection{Inverse Function Compositions}
\label{subsec:inverse-trig-compositions}
%%%%%%%%%%%%%%%%%%%%%%%%%%%%%%%%%%%%%%%%%%%%%%%%%%%%%%%%%%%%

For

\[
x\in[-1,1],
\]

\[
\boxed{
\sin(\arcsin x)
=
x.
}
\]

Likewise,

\[
\boxed{
\cos(\arccos x)
=
x.
}
\]

For every real \(x\),

\[
\boxed{
\tan(\arctan x)
=
x.
}
\]

%%%%%%%%%%%%%%%%%%%%%%%%%%%%%%%%%%%%%%%%%%%%%%%%%%%%%%%%%%%%
\subsection{Reverse Compositions Require Range Restrictions}
\label{subsec:reverse-inverse-trig-compositions}
%%%%%%%%%%%%%%%%%%%%%%%%%%%%%%%%%%%%%%%%%%%%%%%%%%%%%%%%%%%%

In general,

\[
\arcsin(\sin\theta)
\]

does not equal \(\theta\) for every real \(\theta\).

It equals \(\theta\) only when

\[
\boxed{
-\frac{\pi}{2}
\leq
\theta
\leq
\frac{\pi}{2}.
}
\]

Similarly,

\[
\arccos(\cos\theta)
=
\theta
\]

only when

\[
\boxed{
0\leq\theta\leq\pi.
}
\]

%%%%%%%%%%%%%%%%%%%%%%%%%%%%%%%%%%%%%%%%%%%%%%%%%%%%%%%%%%%%
\subsection{Inverse Trig Triangle Method}
\label{subsec:inverse-trig-triangle-method}
%%%%%%%%%%%%%%%%%%%%%%%%%%%%%%%%%%%%%%%%%%%%%%%%%%%%%%%%%%%%

Suppose

\[
\theta
=
\arctan x.
\]

Then

\[
\tan\theta=x.
\]

A right triangle may be chosen with:

\[
\boxed{
\text{opposite}=x,
\qquad
\text{adjacent}=1.
}
\]

Thus the hypotenuse is

\[
\boxed{
\sqrt{1+x^2}.
}
\]

Therefore,

\[
\boxed{
\sin(\arctan x)
=
\frac{x}{\sqrt{1+x^2}},
}
\]

and

\[
\boxed{
\cos(\arctan x)
=
\frac{1}{\sqrt{1+x^2}}.
}
\]

%%%%%%%%%%%%%%%%%%%%%%%%%%%%%%%%%%%%%%%%%%%%%%%%%%%%%%%%%%%%
\subsection{Useful Inverse-Trig Compositions}
\label{subsec:useful-inverse-trig-compositions}
%%%%%%%%%%%%%%%%%%%%%%%%%%%%%%%%%%%%%%%%%%%%%%%%%%%%%%%%%%%%

For appropriate domains,

\[
\boxed{
\sin(\arccos x)
=
\sqrt{1-x^2},
}
\]

because

\[
\arccos x
\in[0,\pi]
\]

and sine is nonnegative there.

Similarly,

\[
\boxed{
\cos(\arcsin x)
=
\sqrt{1-x^2}.
}
\]

%%%%%%%%%%%%%%%%%%%%%%%%%%%%%%%%%%%%%%%%%%%%%%%%%%%%%%%%%%%%
\subsection{Tangent of Arcsine}
\label{subsec:tangent-arcsine}
%%%%%%%%%%%%%%%%%%%%%%%%%%%%%%%%%%%%%%%%%%%%%%%%%%%%%%%%%%%%

For

\[
-1<x<1,
\]

\[
\boxed{
\tan(\arcsin x)
=
\frac{x}{\sqrt{1-x^2}}.
}
\]

%%%%%%%%%%%%%%%%%%%%%%%%%%%%%%%%%%%%%%%%%%%%%%%%%%%%%%%%%%%%
\subsection{Tangent of Arccosine}
\label{subsec:tangent-arccosine}
%%%%%%%%%%%%%%%%%%%%%%%%%%%%%%%%%%%%%%%%%%%%%%%%%%%%%%%%%%%%

For

\[
-1<x<1,
\]

\[
\boxed{
\tan(\arccos x)
=
\frac{\sqrt{1-x^2}}{x}
}
\]

where the sign and domain interpretation must agree with the principal range
and with \(x\neq0\).

%%%%%%%%%%%%%%%%%%%%%%%%%%%%%%%%%%%%%%%%%%%%%%%%%%%%%%%%%%%%
\subsection{Arctangent Complement Identity}
\label{subsec:arctangent-complement}
%%%%%%%%%%%%%%%%%%%%%%%%%%%%%%%%%%%%%%%%%%%%%%%%%%%%%%%%%%%%

For \(x>0\),

\[
\boxed{
\arctan x
+
\arctan\left(\frac1x\right)
=
\frac{\pi}{2}.
}
\]

For \(x<0\), the principal-value relation changes by a sign.

%%%%%%%%%%%%%%%%%%%%%%%%%%%%%%%%%%%%%%%%%%%%%%%%%%%%%%%%%%%%
\subsection{Tangent Addition for Inverse Tangents}
\label{subsec:inverse-tangent-addition}
%%%%%%%%%%%%%%%%%%%%%%%%%%%%%%%%%%%%%%%%%%%%%%%%%%%%%%%%%%%%

If

\[
\alpha=\arctan x,
\qquad
\beta=\arctan y,
\]

then

\[
\boxed{
\tan(\alpha+\beta)
=
\frac{x+y}{1-xy}
}
\]

where the denominator is nonzero.

Recovering \(\alpha+\beta\) from inverse tangent requires quadrant and
principal-value care.

%%%%%%%%%%%%%%%%%%%%%%%%%%%%%%%%%%%%%%%%%%%%%%%%%%%%%%%%%%%%
\subsection{Weierstrass Substitution}
\label{subsec:weierstrass-substitution-reference}
%%%%%%%%%%%%%%%%%%%%%%%%%%%%%%%%%%%%%%%%%%%%%%%%%%%%%%%%%%%%

Set

\[
\boxed{
t
=
\tan\left(\frac{x}{2}\right).
}
\]

Then

\[
\boxed{
\sin x
=
\frac{2t}{1+t^2},
}
\]

\[
\boxed{
\cos x
=
\frac{1-t^2}{1+t^2},
}
\]

and

\[
\boxed{
dx
=
\frac{2}{1+t^2}\,dt.
}
\]

This converts rational expressions in sine and cosine into rational
functions of \(t\).

%%%%%%%%%%%%%%%%%%%%%%%%%%%%%%%%%%%%%%%%%%%%%%%%%%%%%%%%%%%%
\subsection{Pythagorean Parameterization of the Unit Circle}
\label{subsec:unit-circle-rational-parameterization}
%%%%%%%%%%%%%%%%%%%%%%%%%%%%%%%%%%%%%%%%%%%%%%%%%%%%%%%%%%%%

Using

\[
t
=
\tan\left(\frac{\theta}{2}\right),
\]

the unit circle can be parameterized by

\[
\boxed{
x
=
\frac{1-t^2}{1+t^2},
}
\]

\[
\boxed{
y
=
\frac{2t}{1+t^2}.
}
\]

Indeed,

\[
\boxed{
x^2+y^2=1.
}
\]

%%%%%%%%%%%%%%%%%%%%%%%%%%%%%%%%%%%%%%%%%%%%%%%%%%%%%%%%%%%%
\subsection{Phase-Shift Form}
\label{subsec:phase-shift-amplitude-form}
%%%%%%%%%%%%%%%%%%%%%%%%%%%%%%%%%%%%%%%%%%%%%%%%%%%%%%%%%%%%

A sinusoid

\[
A\sin(Bx-C)+D
\]

can be written

\[
\boxed{
A
\sin
\left[
B
\left(
x-\frac{C}{B}
\right)
\right]
+
D
}
\]

when \(B\neq0\).

Thus:

\[
\boxed{
|A|
=
\text{amplitude},
}
\]

\[
\boxed{
\frac{2\pi}{|B|}
=
\text{period},
}
\]

\[
\boxed{
\frac{C}{B}
=
\text{phase shift},
}
\]

and

\[
\boxed{
y=D
=
\text{midline}.
}
\]

%%%%%%%%%%%%%%%%%%%%%%%%%%%%%%%%%%%%%%%%%%%%%%%%%%%%%%%%%%%%
\subsection{Tangent Transformation Form}
\label{subsec:tangent-transformation-reference}
%%%%%%%%%%%%%%%%%%%%%%%%%%%%%%%%%%%%%%%%%%%%%%%%%%%%%%%%%%%%

For

\[
\boxed{
y
=
A\tan(B(x-C))+D,
}
\]

the period is

\[
\boxed{
\frac{\pi}{|B|}.
}
\]

The center point is

\[
\boxed{
(C,D).
}
\]

Vertical asymptotes occur when

\[
B(x-C)
=
\frac{\pi}{2}
+
k\pi.
\]

Therefore,

\[
\boxed{
x
=
C
+
\frac{
\pi/2+k\pi
}{
B
}.
}
\]

%%%%%%%%%%%%%%%%%%%%%%%%%%%%%%%%%%%%%%%%%%%%%%%%%%%%%%%%%%%%
\subsection{Secant and Cosecant Periods}
\label{subsec:sec-csc-periods}
%%%%%%%%%%%%%%%%%%%%%%%%%%%%%%%%%%%%%%%%%%%%%%%%%%%%%%%%%%%%

For

\[
y=A\sec(Bx-C)+D
\]

and

\[
y=A\csc(Bx-C)+D,
\]

the period is

\[
\boxed{
\frac{2\pi}{|B|}.
}
\]

Their vertical asymptotes occur where their reciprocal denominator vanishes.

%%%%%%%%%%%%%%%%%%%%%%%%%%%%%%%%%%%%%%%%%%%%%%%%%%%%%%%%%%%%
\subsection{Identity Verification Strategy}
\label{subsec:identity-verification-strategy}
%%%%%%%%%%%%%%%%%%%%%%%%%%%%%%%%%%%%%%%%%%%%%%%%%%%%%%%%%%%%

When verifying a trigonometric identity:

\begin{enumerate}
    \item usually work on only one side at a time;
    \item convert secant, cosecant, tangent, and cotangent to sine and cosine
          when helpful;
    \item use Pythagorean identities strategically;
    \item factor before canceling;
    \item combine fractions using a common denominator;
    \item use conjugates when expressions such as \(1\pm\sin x\) occur;
    \item apply sum, difference, or double-angle identities when patterns
          appear;
    \item preserve domain restrictions;
    \item stop once the transformed side matches the other side.
\end{enumerate}

%%%%%%%%%%%%%%%%%%%%%%%%%%%%%%%%%%%%%%%%%%%%%%%%%%%%%%%%%%%%
\subsection{Worked Example: Verify an Identity}
\label{subsec:worked-verify-identity}
%%%%%%%%%%%%%%%%%%%%%%%%%%%%%%%%%%%%%%%%%%%%%%%%%%%%%%%%%%%%

\begin{workedexamplebox}{Verify a Reciprocal Identity}

Verify

\[
\frac{1-\cos^2x}{\sin x}
=
\sin x.
\]

Start with the left side.

Using

\[
1-\cos^2x
=
\sin^2x,
\]

we obtain

\[
\frac{\sin^2x}{\sin x}.
\]

For \(\sin x\neq0\),

\[
\frac{\sin^2x}{\sin x}
=
\sin x.
\]

Thus

\begin{answerbox}
\[
\boxed{
\frac{1-\cos^2x}{\sin x}
=
\sin x
}
\]

where the original left side is defined.
\end{answerbox}

\end{workedexamplebox}

%%%%%%%%%%%%%%%%%%%%%%%%%%%%%%%%%%%%%%%%%%%%%%%%%%%%%%%%%%%%
\subsection{Worked Example: Identity With Secant}
\label{subsec:worked-verify-secant}
%%%%%%%%%%%%%%%%%%%%%%%%%%%%%%%%%%%%%%%%%%%%%%%%%%%%%%%%%%%%

\begin{workedexamplebox}{Simplify \(\frac{\sec x-\cos x}{\tan x}\)}

Convert secant and tangent:

\[
\frac{
\frac1{\cos x}-\cos x
}{
\frac{\sin x}{\cos x}
}.
\]

Simplify the numerator:

\[
\frac{
\frac{
1-\cos^2x
}{
\cos x
}
}{
\frac{\sin x}{\cos x}
}.
\]

Use

\[
1-\cos^2x
=
\sin^2x.
\]

Then

\[
\frac{
\sin^2x/\cos x
}{
\sin x/\cos x
}
=
\sin x.
\]

Therefore,

\begin{answerbox}
\[
\boxed{
\frac{\sec x-\cos x}{\tan x}
=
\sin x
}
\]

where the original expression is defined.
\end{answerbox}

\end{workedexamplebox}

%%%%%%%%%%%%%%%%%%%%%%%%%%%%%%%%%%%%%%%%%%%%%%%%%%%%%%%%%%%%
\subsection{Worked Example: Conjugate Technique}
\label{subsec:worked-trig-conjugate-technique}
%%%%%%%%%%%%%%%%%%%%%%%%%%%%%%%%%%%%%%%%%%%%%%%%%%%%%%%%%%%%

\begin{workedexamplebox}{Simplify \(\frac{1-\cos x}{\sin x}\)}

Multiply numerator and denominator by

\[
1+\cos x.
\]

Then

\[
\frac{1-\cos x}{\sin x}
\cdot
\frac{1+\cos x}{1+\cos x}
=
\frac{
1-\cos^2x
}{
\sin x(1+\cos x)
}.
\]

Using

\[
1-\cos^2x
=
\sin^2x,
\]

we get

\[
\frac{
\sin^2x
}{
\sin x(1+\cos x)
}
=
\frac{
\sin x
}{
1+\cos x
}.
\]

Thus

\begin{answerbox}
\[
\boxed{
\frac{1-\cos x}{\sin x}
=
\frac{\sin x}{1+\cos x}
}
\]

where both original expressions are defined.
\end{answerbox}

\end{workedexamplebox}

%%%%%%%%%%%%%%%%%%%%%%%%%%%%%%%%%%%%%%%%%%%%%%%%%%%%%%%%%%%%
\subsection{Small-Angle Identities and Approximations}
\label{subsec:small-angle-reference}
%%%%%%%%%%%%%%%%%%%%%%%%%%%%%%%%%%%%%%%%%%%%%%%%%%%%%%%%%%%%

These are not exact identities, but important approximations.

As

\[
x\to0,
\]

\[
\boxed{
\sin x
\sim
x,
}
\]

\[
\boxed{
\tan x
\sim
x,
}
\]

and

\[
\boxed{
1-\cos x
\sim
\frac{x^2}{2}.
}
\]

These require \(x\) measured in radians.

%%%%%%%%%%%%%%%%%%%%%%%%%%%%%%%%%%%%%%%%%%%%%%%%%%%%%%%%%%%%
\subsection{Fundamental Trigonometric Limits}
\label{subsec:fundamental-trig-limits-reference}
%%%%%%%%%%%%%%%%%%%%%%%%%%%%%%%%%%%%%%%%%%%%%%%%%%%%%%%%%%%%

In radians,

\[
\boxed{
\lim_{x\to0}
\frac{\sin x}{x}
=
1,
}
\]

and

\[
\boxed{
\lim_{x\to0}
\frac{1-\cos x}{x^2}
=
\frac12.
}
\]

These connect trigonometric identities with calculus.

%%%%%%%%%%%%%%%%%%%%%%%%%%%%%%%%%%%%%%%%%%%%%%%%%%%%%%%%%%%%
\subsection{Derivative Identities}
\label{subsec:trig-derivative-reference}
%%%%%%%%%%%%%%%%%%%%%%%%%%%%%%%%%%%%%%%%%%%%%%%%%%%%%%%%%%%%

The standard derivatives are

\[
\boxed{
\frac{d}{dx}\sin x
=
\cos x,
}
\]

\[
\boxed{
\frac{d}{dx}\cos x
=
-\sin x,
}
\]

\[
\boxed{
\frac{d}{dx}\tan x
=
\sec^2x,
}
\]

\[
\boxed{
\frac{d}{dx}\cot x
=
-\csc^2x,
}
\]

\[
\boxed{
\frac{d}{dx}\sec x
=
\sec x\tan x,
}
\]

and

\[
\boxed{
\frac{d}{dx}\csc x
=
-\csc x\cot x.
}
\]

%%%%%%%%%%%%%%%%%%%%%%%%%%%%%%%%%%%%%%%%%%%%%%%%%%%%%%%%%%%%
\subsection{Integral Connections}
\label{subsec:trig-integral-reference}
%%%%%%%%%%%%%%%%%%%%%%%%%%%%%%%%%%%%%%%%%%%%%%%%%%%%%%%%%%%%

Important antiderivatives include

\[
\boxed{
\int\sin x\,dx
=
-\cos x+C,
}
\]

\[
\boxed{
\int\cos x\,dx
=
\sin x+C,
}
\]

\[
\boxed{
\int\sec^2x\,dx
=
\tan x+C,
}
\]

\[
\boxed{
\int\csc^2x\,dx
=
-\cot x+C.
}
\]

%%%%%%%%%%%%%%%%%%%%%%%%%%%%%%%%%%%%%%%%%%%%%%%%%%%%%%%%%%%%
\subsection{Power-Reduction in Integration}
\label{subsec:power-reduction-integration}
%%%%%%%%%%%%%%%%%%%%%%%%%%%%%%%%%%%%%%%%%%%%%%%%%%%%%%%%%%%%

For integrals involving even powers,

\[
\boxed{
\sin^2x
=
\frac{1-\cos2x}{2},
}
\]

and

\[
\boxed{
\cos^2x
=
\frac{1+\cos2x}{2}
}
\]

are especially useful.

For example,

\[
\int\sin^2x\,dx
=
\frac12
\int
(1-\cos2x)\,dx.
\]

%%%%%%%%%%%%%%%%%%%%%%%%%%%%%%%%%%%%%%%%%%%%%%%%%%%%%%%%%%%%
\subsection{Product-to-Sum in Integration}
\label{subsec:product-to-sum-integration}
%%%%%%%%%%%%%%%%%%%%%%%%%%%%%%%%%%%%%%%%%%%%%%%%%%%%%%%%%%%%

Products such as

\[
\sin(mx)\cos(nx)
\]

can be converted into sums:

\[
\boxed{
\sin(mx)\cos(nx)
=
\frac12
\left[
\sin((m+n)x)
+
\sin((m-n)x)
\right].
}
\]

This is useful in Fourier analysis and integration.

%%%%%%%%%%%%%%%%%%%%%%%%%%%%%%%%%%%%%%%%%%%%%%%%%%%%%%%%%%%%
\subsection{Orthogonality of Sine and Cosine}
\label{subsec:sine-cosine-orthogonality}
%%%%%%%%%%%%%%%%%%%%%%%%%%%%%%%%%%%%%%%%%%%%%%%%%%%%%%%%%%%%

For positive integers \(m,n\),

\[
\boxed{
\int_{-\pi}^{\pi}
\sin(mx)\sin(nx)\,dx
=
0
\quad
\text{if }m\neq n,
}
\]

and

\[
\boxed{
\int_{-\pi}^{\pi}
\cos(mx)\cos(nx)\,dx
=
0
\quad
\text{if }m\neq n.
}
\]

Also,

\[
\boxed{
\int_{-\pi}^{\pi}
\sin(mx)\cos(nx)\,dx
=
0.
}
\]

These identities are foundational in Fourier series.

%%%%%%%%%%%%%%%%%%%%%%%%%%%%%%%%%%%%%%%%%%%%%%%%%%%%%%%%%%%%
\subsection{Norms of Fourier Basis Functions}
\label{subsec:fourier-basis-norms}
%%%%%%%%%%%%%%%%%%%%%%%%%%%%%%%%%%%%%%%%%%%%%%%%%%%%%%%%%%%%

For positive integer \(n\),

\[
\boxed{
\int_{-\pi}^{\pi}
\sin^2(nx)\,dx
=
\pi,
}
\]

and

\[
\boxed{
\int_{-\pi}^{\pi}
\cos^2(nx)\,dx
=
\pi.
}
\]

Also,

\[
\boxed{
\int_{-\pi}^{\pi}1\,dx
=
2\pi.
}
\]

%%%%%%%%%%%%%%%%%%%%%%%%%%%%%%%%%%%%%%%%%%%%%%%%%%%%%%%%%%%%
\subsection{Harmonic Addition}
\label{subsec:harmonic-addition-reference}
%%%%%%%%%%%%%%%%%%%%%%%%%%%%%%%%%%%%%%%%%%%%%%%%%%%%%%%%%%%%

Two sinusoidal signals with the same frequency can be combined into a single
sinusoid.

For example,

\[
a\cos(\omega t)
+
b\sin(\omega t)
=
R
\cos(\omega t-\phi),
\]

where

\[
\boxed{
R
=
\sqrt{a^2+b^2}.
}
\]

This is important in oscillation and signal models.

%%%%%%%%%%%%%%%%%%%%%%%%%%%%%%%%%%%%%%%%%%%%%%%%%%%%%%%%%%%%
\subsection{Trigonometric Identities in Differential Equations}
\label{subsec:trig-identities-differential-equations}
%%%%%%%%%%%%%%%%%%%%%%%%%%%%%%%%%%%%%%%%%%%%%%%%%%%%%%%%%%%%

Solutions of the harmonic oscillator

\[
y''+\omega^2y=0
\]

can be written

\[
\boxed{
y
=
A\cos(\omega t)
+
B\sin(\omega t).
}
\]

Using auxiliary-angle form,

\[
\boxed{
y
=
R\cos(\omega t-\phi).
}
\]

Thus two integration constants can be reinterpreted as amplitude and phase.

%%%%%%%%%%%%%%%%%%%%%%%%%%%%%%%%%%%%%%%%%%%%%%%%%%%%%%%%%%%%
\subsection{Trigonometric Identities in Complex Analysis}
\label{subsec:trig-identities-complex-analysis}
%%%%%%%%%%%%%%%%%%%%%%%%%%%%%%%%%%%%%%%%%%%%%%%%%%%%%%%%%%%%

Euler's formula gives

\[
e^{iz}
=
\cos z+i\sin z
\]

for complex \(z\).

Thus one may define

\[
\boxed{
\cos z
=
\frac{
e^{iz}+e^{-iz}
}{2},
}
\]

and

\[
\boxed{
\sin z
=
\frac{
e^{iz}-e^{-iz}
}{
2i
}.
}
\]

These extend trigonometric functions to the complex plane.

%%%%%%%%%%%%%%%%%%%%%%%%%%%%%%%%%%%%%%%%%%%%%%%%%%%%%%%%%%%%
\subsection{Trigonometric and Hyperbolic Connections}
\label{subsec:trig-hyperbolic-connections}
%%%%%%%%%%%%%%%%%%%%%%%%%%%%%%%%%%%%%%%%%%%%%%%%%%%%%%%%%%%%

For real \(x\),

\[
\boxed{
\cos(ix)
=
\cosh x,
}
\]

and

\[
\boxed{
\sin(ix)
=
i\sinh x.
}
\]

Equivalently,

\[
\boxed{
\cosh x
=
\cos(ix),
}
\]

and

\[
\boxed{
\sinh x
=
-i\sin(ix).
}
\]

%%%%%%%%%%%%%%%%%%%%%%%%%%%%%%%%%%%%%%%%%%%%%%%%%%%%%%%%%%%%
\subsection{Common Errors}
\label{subsec:trigonometric-identities-common-errors}
%%%%%%%%%%%%%%%%%%%%%%%%%%%%%%%%%%%%%%%%%%%%%%%%%%%%%%%%%%%%

\subsubsection{Confusing \(\sin^2x\) With \(\sin(x^2)\)}

The notation

\[
\boxed{
\sin^2x
}
\]

means

\[
\boxed{
(\sin x)^2.
}
\]

It does not mean

\[
\sin(x^2).
\]

%%%%%%%%%%%%%%%%%%%%%%%%%%%%%%%%%%%%%%%%%%%%%%%%%%%%%%%%%%%%
\subsubsection{Confusing \(\sin^{-1}x\) With Reciprocal}

In inverse-function notation,

\[
\boxed{
\sin^{-1}x
=
\arcsin x.
}
\]

But reciprocal sine is

\[
\boxed{
\csc x
=
\frac1{\sin x}.
}
\]

These are different.

%%%%%%%%%%%%%%%%%%%%%%%%%%%%%%%%%%%%%%%%%%%%%%%%%%%%%%%%%%%%
\subsubsection{Incorrect Sine Addition}
\label{subsec:false-sine-addition}
%%%%%%%%%%%%%%%%%%%%%%%%%%%%%%%%%%%%%%%%%%%%%%%%%%%%%%%%%%%%

In general,

\[
\boxed{
\sin(\alpha+\beta)
\neq
\sin\alpha+\sin\beta.
}
\]

The correct identity is

\[
\boxed{
\sin(\alpha+\beta)
=
\sin\alpha\cos\beta
+
\cos\alpha\sin\beta.
}
\]

%%%%%%%%%%%%%%%%%%%%%%%%%%%%%%%%%%%%%%%%%%%%%%%%%%%%%%%%%%%%
\subsubsection{Incorrect Cosine Addition}
%%%%%%%%%%%%%%%%%%%%%%%%%%%%%%%%%%%%%%%%%%%%%%%%%%%%%%%%%%%%

In general,

\[
\boxed{
\cos(\alpha+\beta)
\neq
\cos\alpha+\cos\beta.
}
\]

The correct formula is

\[
\boxed{
\cos(\alpha+\beta)
=
\cos\alpha\cos\beta
-
\sin\alpha\sin\beta.
}
\]

%%%%%%%%%%%%%%%%%%%%%%%%%%%%%%%%%%%%%%%%%%%%%%%%%%%%%%%%%%%%
\subsubsection{Forgetting the Sign in Cosine Sum and Difference}
%%%%%%%%%%%%%%%%%%%%%%%%%%%%%%%%%%%%%%%%%%%%%%%%%%%%%%%%%%%%

Remember:

\[
\boxed{
\cos(\alpha+\beta)
=
\cos\alpha\cos\beta
-
\sin\alpha\sin\beta,
}
\]

while

\[
\boxed{
\cos(\alpha-\beta)
=
\cos\alpha\cos\beta
+
\sin\alpha\sin\beta.
}
\]

The sign inside the angle is reversed in the middle term.

%%%%%%%%%%%%%%%%%%%%%%%%%%%%%%%%%%%%%%%%%%%%%%%%%%%%%%%%%%%%
\subsubsection{Forgetting the Factor of Two}
%%%%%%%%%%%%%%%%%%%%%%%%%%%%%%%%%%%%%%%%%%%%%%%%%%%%%%%%%%%%

The sine double-angle identity is

\[
\boxed{
\sin2x
=
2\sin x\cos x,
}
\]

not simply

\[
\sin x\cos x.
\]

%%%%%%%%%%%%%%%%%%%%%%%%%%%%%%%%%%%%%%%%%%%%%%%%%%%%%%%%%%%%
\subsubsection{Incorrect Square Root From Half-Angle}
%%%%%%%%%%%%%%%%%%%%%%%%%%%%%%%%%%%%%%%%%%%%%%%%%%%%%%%%%%%%

From

\[
\sin^2
\left(
\frac{x}{2}
\right)
=
\frac{1-\cos x}{2},
\]

we get

\[
\boxed{
\sin
\left(
\frac{x}{2}
\right)
=
\pm
\sqrt{
\frac{1-\cos x}{2}
}.
}
\]

The sign depends on the quadrant.

%%%%%%%%%%%%%%%%%%%%%%%%%%%%%%%%%%%%%%%%%%%%%%%%%%%%%%%%%%%%
\subsubsection{Ignoring Trigonometric Domains}
%%%%%%%%%%%%%%%%%%%%%%%%%%%%%%%%%%%%%%%%%%%%%%%%%%%%%%%%%%%%

For example,

\[
\tan x
\]

is undefined when

\[
\boxed{
\cos x=0.
}
\]

Similarly,

\[
\sec x
\]

is undefined at the same points.

%%%%%%%%%%%%%%%%%%%%%%%%%%%%%%%%%%%%%%%%%%%%%%%%%%%%%%%%%%%%
\subsubsection{Canceling Terms Instead of Factors}
%%%%%%%%%%%%%%%%%%%%%%%%%%%%%%%%%%%%%%%%%%%%%%%%%%%%%%%%%%%%

In general,

\[
\frac{
\sin x+\cos x
}{
\sin x
}
\]

does not simplify by canceling the \(\sin x\) term.

Instead,

\[
\boxed{
\frac{
\sin x+\cos x
}{
\sin x
}
=
1+\cot x.
}
\]

%%%%%%%%%%%%%%%%%%%%%%%%%%%%%%%%%%%%%%%%%%%%%%%%%%%%%%%%%%%%
\subsubsection{Dividing by a Trigonometric Function Without Checking Zero}
%%%%%%%%%%%%%%%%%%%%%%%%%%%%%%%%%%%%%%%%%%%%%%%%%%%%%%%%%%%%

From

\[
\sin x\,A
=
\sin x\,B,
\]

one may divide by \(\sin x\) only when

\[
\boxed{
\sin x\neq0.
}
\]

Otherwise possible solutions can be lost.

%%%%%%%%%%%%%%%%%%%%%%%%%%%%%%%%%%%%%%%%%%%%%%%%%%%%%%%%%%%%
\subsubsection{Using Degree-Based Angles in Calculus Limits}
%%%%%%%%%%%%%%%%%%%%%%%%%%%%%%%%%%%%%%%%%%%%%%%%%%%%%%%%%%%%

The fundamental limit

\[
\boxed{
\lim_{x\to0}
\frac{\sin x}{x}
=
1
}
\]

requires radians.

%%%%%%%%%%%%%%%%%%%%%%%%%%%%%%%%%%%%%%%%%%%%%%%%%%%%%%%%%%%%
\subsubsection{Assuming Inverse-Trig Composition Always Cancels}
%%%%%%%%%%%%%%%%%%%%%%%%%%%%%%%%%%%%%%%%%%%%%%%%%%%%%%%%%%%%

The identity

\[
\arcsin(\sin x)=x
\]

does not hold for every real \(x\).

Principal-value ranges must be respected.

%%%%%%%%%%%%%%%%%%%%%%%%%%%%%%%%%%%%%%%%%%%%%%%%%%%%%%%%%%%%
\subsubsection{Ignoring Quadrants in Inverse Tangent Sums}
%%%%%%%%%%%%%%%%%%%%%%%%%%%%%%%%%%%%%%%%%%%%%%%%%%%%%%%%%%%%

Even if

\[
\tan\theta
=
z,
\]

the equation

\[
\theta=\arctan z
\]

determines only the principal value.

Additional multiples of \(\pi\) may be needed.

%%%%%%%%%%%%%%%%%%%%%%%%%%%%%%%%%%%%%%%%%%%%%%%%%%%%%%%%%%%%
\subsubsection{Confusing Identity Verification With Equation Solving}
%%%%%%%%%%%%%%%%%%%%%%%%%%%%%%%%%%%%%%%%%%%%%%%%%%%%%%%%%%%%

When verifying an identity, the goal is to show two expressions agree on
their common domain.

When solving an equation, the goal is to find specific values of the
variable.

These are different tasks.

%%%%%%%%%%%%%%%%%%%%%%%%%%%%%%%%%%%%%%%%%%%%%%%%%%%%%%%%%%%%
\subsection{Quick Reciprocal and Quotient Table}
\label{subsec:quick-reciprocal-quotient-trig}
%%%%%%%%%%%%%%%%%%%%%%%%%%%%%%%%%%%%%%%%%%%%%%%%%%%%%%%%%%%%

\[
\boxed{
\csc x
=
\frac1{\sin x}
}
\]

\[
\boxed{
\sec x
=
\frac1{\cos x}
}
\]

\[
\boxed{
\cot x
=
\frac1{\tan x}
}
\]

\[
\boxed{
\tan x
=
\frac{\sin x}{\cos x}
}
\]

\[
\boxed{
\cot x
=
\frac{\cos x}{\sin x}
}
\]

%%%%%%%%%%%%%%%%%%%%%%%%%%%%%%%%%%%%%%%%%%%%%%%%%%%%%%%%%%%%
\subsection{Quick Pythagorean Table}
\label{subsec:quick-pythagorean-trig}
%%%%%%%%%%%%%%%%%%%%%%%%%%%%%%%%%%%%%%%%%%%%%%%%%%%%%%%%%%%%

\[
\boxed{
\sin^2x+\cos^2x=1
}
\]

\[
\boxed{
1+\tan^2x=\sec^2x
}
\]

\[
\boxed{
1+\cot^2x=\csc^2x
}
\]

%%%%%%%%%%%%%%%%%%%%%%%%%%%%%%%%%%%%%%%%%%%%%%%%%%%%%%%%%%%%
\subsection{Quick Sum and Difference Table}
\label{subsec:quick-sum-difference-trig}
%%%%%%%%%%%%%%%%%%%%%%%%%%%%%%%%%%%%%%%%%%%%%%%%%%%%%%%%%%%%

\[
\boxed{
\sin(\alpha+\beta)
=
\sin\alpha\cos\beta
+
\cos\alpha\sin\beta
}
\]

\[
\boxed{
\sin(\alpha-\beta)
=
\sin\alpha\cos\beta
-
\cos\alpha\sin\beta
}
\]

\[
\boxed{
\cos(\alpha+\beta)
=
\cos\alpha\cos\beta
-
\sin\alpha\sin\beta
}
\]

\[
\boxed{
\cos(\alpha-\beta)
=
\cos\alpha\cos\beta
+
\sin\alpha\sin\beta
}
\]

\[
\boxed{
\tan(\alpha+\beta)
=
\frac{
\tan\alpha+\tan\beta
}{
1-\tan\alpha\tan\beta
}
}
\]

\[
\boxed{
\tan(\alpha-\beta)
=
\frac{
\tan\alpha-\tan\beta
}{
1+\tan\alpha\tan\beta
}
}
\]

%%%%%%%%%%%%%%%%%%%%%%%%%%%%%%%%%%%%%%%%%%%%%%%%%%%%%%%%%%%%
\subsection{Quick Double-Angle Table}
\label{subsec:quick-double-angle-trig}
%%%%%%%%%%%%%%%%%%%%%%%%%%%%%%%%%%%%%%%%%%%%%%%%%%%%%%%%%%%%

\[
\boxed{
\sin2x
=
2\sin x\cos x
}
\]

\[
\boxed{
\cos2x
=
\cos^2x-\sin^2x
}
\]

\[
\boxed{
\cos2x
=
2\cos^2x-1
}
\]

\[
\boxed{
\cos2x
=
1-2\sin^2x
}
\]

\[
\boxed{
\tan2x
=
\frac{
2\tan x
}{
1-\tan^2x
}
}
\]

%%%%%%%%%%%%%%%%%%%%%%%%%%%%%%%%%%%%%%%%%%%%%%%%%%%%%%%%%%%%
\subsection{Quick Half-Angle Table}
\label{subsec:quick-half-angle-trig}
%%%%%%%%%%%%%%%%%%%%%%%%%%%%%%%%%%%%%%%%%%%%%%%%%%%%%%%%%%%%

\[
\boxed{
\sin^2\left(\frac{x}{2}\right)
=
\frac{1-\cos x}{2}
}
\]

\[
\boxed{
\cos^2\left(\frac{x}{2}\right)
=
\frac{1+\cos x}{2}
}
\]

\[
\boxed{
\tan\left(\frac{x}{2}\right)
=
\frac{\sin x}{1+\cos x}
=
\frac{1-\cos x}{\sin x}
}
\]

%%%%%%%%%%%%%%%%%%%%%%%%%%%%%%%%%%%%%%%%%%%%%%%%%%%%%%%%%%%%
\subsection{Quick Product-to-Sum Table}
\label{subsec:quick-product-to-sum-trig}
%%%%%%%%%%%%%%%%%%%%%%%%%%%%%%%%%%%%%%%%%%%%%%%%%%%%%%%%%%%%

\[
\boxed{
\sin A\sin B
=
\frac12
[
\cos(A-B)-\cos(A+B)
]
}
\]

\[
\boxed{
\cos A\cos B
=
\frac12
[
\cos(A-B)+\cos(A+B)
]
}
\]

\[
\boxed{
\sin A\cos B
=
\frac12
[
\sin(A+B)+\sin(A-B)
]
}
\]

\[
\boxed{
\cos A\sin B
=
\frac12
[
\sin(A+B)-\sin(A-B)
]
}
\]

%%%%%%%%%%%%%%%%%%%%%%%%%%%%%%%%%%%%%%%%%%%%%%%%%%%%%%%%%%%%
\subsection{Quick Sum-to-Product Table}
\label{subsec:quick-sum-to-product-trig}
%%%%%%%%%%%%%%%%%%%%%%%%%%%%%%%%%%%%%%%%%%%%%%%%%%%%%%%%%%%%

\[
\boxed{
\sin A+\sin B
=
2
\sin\left(\frac{A+B}{2}\right)
\cos\left(\frac{A-B}{2}\right)
}
\]

\[
\boxed{
\sin A-\sin B
=
2
\cos\left(\frac{A+B}{2}\right)
\sin\left(\frac{A-B}{2}\right)
}
\]

\[
\boxed{
\cos A+\cos B
=
2
\cos\left(\frac{A+B}{2}\right)
\cos\left(\frac{A-B}{2}\right)
}
\]

\[
\boxed{
\cos A-\cos B
=
-2
\sin\left(\frac{A+B}{2}\right)
\sin\left(\frac{A-B}{2}\right)
}
\]

%%%%%%%%%%%%%%%%%%%%%%%%%%%%%%%%%%%%%%%%%%%%%%%%%%%%%%%%%%%%
\subsection{Quick Special-Angle Table}
\label{subsec:quick-special-angle-table}
%%%%%%%%%%%%%%%%%%%%%%%%%%%%%%%%%%%%%%%%%%%%%%%%%%%%%%%%%%%%

\begin{center}
\begin{tabular}{c|c|c|c}
\textbf{Angle}
&
\(\sin\theta\)
&
\(\cos\theta\)
&
\(\tan\theta\)
\\ \hline

\(0\)
&
\(0\)
&
\(1\)
&
\(0\)
\\

\(\frac{\pi}{6}\)
&
\(\frac12\)
&
\(\frac{\sqrt3}{2}\)
&
\(\frac{\sqrt3}{3}\)
\\

\(\frac{\pi}{4}\)
&
\(\frac{\sqrt2}{2}\)
&
\(\frac{\sqrt2}{2}\)
&
\(1\)
\\

\(\frac{\pi}{3}\)
&
\(\frac{\sqrt3}{2}\)
&
\(\frac12\)
&
\(\sqrt3\)
\\

\(\frac{\pi}{2}\)
&
\(1\)
&
\(0\)
&
undefined
\end{tabular}
\end{center}

%%%%%%%%%%%%%%%%%%%%%%%%%%%%%%%%%%%%%%%%%%%%%%%%%%%%%%%%%%%%
\subsection{Exercises}
\label{subsec:trigonometric-identities-exercises}
%%%%%%%%%%%%%%%%%%%%%%%%%%%%%%%%%%%%%%%%%%%%%%%%%%%%%%%%%%%%

\begin{exercise}[1]
State the reciprocal identity for \(\sec x\).
\end{exercise}

\begin{exercise}[1]
State the quotient identity for \(\tan x\).
\end{exercise}

\begin{exercise}[1]
State all three Pythagorean identities.
\end{exercise}

\begin{exercise}[1]
State whether sine is even or odd.
\end{exercise}

\begin{exercise}[1]
State whether cosine is even or odd.
\end{exercise}

\begin{exercise}[1]
State the period of tangent.
\end{exercise}

\begin{exercise}[1]
State the period of cosine.
\end{exercise}

\begin{exercise}[2]
Simplify

\[
\frac{\sin x}{\cos x}.
\]
\end{exercise}

\begin{exercise}[2]
Simplify

\[
\frac{1}{\sec x}.
\]
\end{exercise}

\begin{exercise}[2]
Simplify

\[
1-\sin^2x.
\]
\end{exercise}

\begin{exercise}[2]
Simplify

\[
\sec^2x-\tan^2x.
\]
\end{exercise}

\begin{exercise}[2]
Evaluate

\[
\sin(-\pi/6).
\]
\end{exercise}

\begin{exercise}[2]
Evaluate

\[
\cos(-\pi/3).
\]
\end{exercise}

\begin{exercise}[2]
Evaluate

\[
\tan(\pi+\pi/4).
\]
\end{exercise}

\begin{exercise}[2]
Use a cofunction identity to rewrite

\[
\sin\left(\frac{\pi}{2}-x\right).
\]
\end{exercise}

\begin{exercise}[2]
Find the exact value of

\[
\sin75^\circ.
\]
\end{exercise}

\begin{exercise}[2]
Find the exact value of

\[
\cos15^\circ.
\]
\end{exercise}

\begin{exercise}[2]
Find the exact value of

\[
\tan75^\circ.
\]
\end{exercise}

\begin{exercise}[2]
Expand

\[
\sin(2x+y).
\]
\end{exercise}

\begin{exercise}[2]
Expand

\[
\cos(3x-y).
\]
\end{exercise}

\begin{exercise}[2]
Rewrite

\[
\sin2x
\]

using only \(\sin x\) and \(\cos x\).
\end{exercise}

\begin{exercise}[2]
Rewrite

\[
\cos2x
\]

in three equivalent forms.
\end{exercise}

\begin{exercise}[2]
Rewrite

\[
\sin^2x
\]

using a double-angle expression.
\end{exercise}

\begin{exercise}[2]
Rewrite

\[
\cos^2x
\]

using a double-angle expression.
\end{exercise}

\begin{exercise}[2]
Find

\[
\sin\left(\frac{x}{2}\right)
\]

given \(\cos x\), with the appropriate quadrant specified.
\end{exercise}

\begin{exercise}[2]
Find

\[
\cos\left(\frac{x}{2}\right)
\]

given \(\cos x\), with quadrant information.
\end{exercise}

\begin{exercise}[2]
Rewrite

\[
\sin3x
\]

using powers of \(\sin x\).
\end{exercise}

\begin{exercise}[2]
Rewrite

\[
\cos3x
\]

using powers of \(\cos x\).
\end{exercise}

\begin{exercise}[2]
Convert

\[
\sin4x\sin2x
\]

to a sum.
\end{exercise}

\begin{exercise}[2]
Convert

\[
\cos5x\cos x
\]

to a sum.
\end{exercise}

\begin{exercise}[2]
Convert

\[
\sin7x\cos3x
\]

to a sum.
\end{exercise}

\begin{exercise}[2]
Convert

\[
\sin5x+\sin x
\]

to a product.
\end{exercise}

\begin{exercise}[2]
Convert

\[
\cos4x+\cos2x
\]

to a product.
\end{exercise}

\begin{exercise}[2]
Simplify

\[
\frac{
1-\cos^2x
}{
\sin x
}.
\]
\end{exercise}

\begin{exercise}[2]
Simplify

\[
\frac{
\sec x-\cos x
}{
\tan x
}.
\]
\end{exercise}

\begin{exercise}[2]
Simplify

\[
\frac{
1-\cos x
}{
\sin x
}
\]

using a conjugate.
\end{exercise}

\begin{exercise}[3]
Verify

\[
\frac{
\sin x
}{
1+\cos x
}
=
\frac{
1-\cos x
}{
\sin x
}.
\]
\end{exercise}

\begin{exercise}[3]
Verify

\[
\frac{
1+\tan^2x
}{
\sec^2x
}
=
1.
\]
\end{exercise}

\begin{exercise}[3]
Verify

\[
(\sec x-\tan x)
(\sec x+\tan x)
=
1.
\]
\end{exercise}

\begin{exercise}[3]
Verify

\[
\frac{
1-\sin x
}{
\cos x
}
=
\frac{
\cos x
}{
1+\sin x
}.
\]
\end{exercise}

\begin{exercise}[3]
Rewrite

\[
5\cos x+12\sin x
\]

as a single sinusoid.
\end{exercise}

\begin{exercise}[3]
Find the maximum possible value of

\[
7\cos x-24\sin x.
\]
\end{exercise}

\begin{exercise}[3]
Derive the sine double-angle identity from the sine sum formula.
\end{exercise}

\begin{exercise}[3]
Derive all three cosine double-angle forms.
\end{exercise}

\begin{exercise}[3]
Derive the tangent double-angle identity.
\end{exercise}

\begin{exercise}[3]
Derive the power-reduction identities.
\end{exercise}

\begin{exercise}[3]
Derive the product-to-sum formulas from sum and difference identities.
\end{exercise}

\begin{exercise}[3]
Derive the sum-to-product formulas.
\end{exercise}

\begin{exercise}[3]
Use De Moivre's formula to derive

\[
\cos3x
=
4\cos^3x-3\cos x.
\]
\end{exercise}

\begin{exercise}[3]
Use De Moivre's formula to derive

\[
\sin3x
=
3\sin x-4\sin^3x.
\]
\end{exercise}

\begin{exercise}[2]
Evaluate

\[
\sin(\arcsin(3/5)).
\]
\end{exercise}

\begin{exercise}[2]
Evaluate

\[
\cos(\arccos(-1/2)).
\]
\end{exercise}

\begin{exercise}[2]
Evaluate

\[
\tan(\arctan 7).
\]
\end{exercise}

\begin{exercise}[2]
Simplify

\[
\sin(\arccos x).
\]
\end{exercise}

\begin{exercise}[2]
Simplify

\[
\cos(\arcsin x).
\]
\end{exercise}

\begin{exercise}[2]
Simplify

\[
\sin(\arctan x).
\]
\end{exercise}

\begin{exercise}[2]
Simplify

\[
\cos(\arctan x).
\]
\end{exercise}

\begin{exercise}[3]
Explain why

\[
\arcsin(\sin(3\pi/4))
\neq
3\pi/4.
\]
\end{exercise}

\begin{exercise}[3]
Evaluate

\[
\arccos(\cos(5\pi/3)).
\]
\end{exercise}

\begin{exercise}[3]
Prove for \(x>0\),

\[
\arctan x
+
\arctan(1/x)
=
\pi/2.
\]
\end{exercise}

\begin{exercise}[3]
Derive the Weierstrass substitutions

\[
\sin x
=
\frac{2t}{1+t^2}
\]

and

\[
\cos x
=
\frac{1-t^2}{1+t^2}.
\]
\end{exercise}

\begin{exercise}[3]
Verify that the rational parameterization

\[
x
=
\frac{1-t^2}{1+t^2},
\qquad
y
=
\frac{2t}{1+t^2}
\]

satisfies

\[
x^2+y^2=1.
\]
\end{exercise}

\begin{exercise}[2]
For

\[
y=3\sin(2(x-\pi/4))+1,
\]

identify amplitude, period, phase shift, and midline.
\end{exercise}

\begin{exercise}[2]
For

\[
y=-2\tan(3(x-\pi/6))+4,
\]

identify period, center point, and vertical asymptotes.
\end{exercise}

\begin{exercise}[3]
Use

\[
\sin x\sim x
\]

to evaluate a basic trigonometric limit.
\end{exercise}

\begin{exercise}[3]
Evaluate

\[
\lim_{x\to0}
\frac{\sin5x}{x}.
\]
\end{exercise}

\begin{exercise}[3]
Evaluate

\[
\lim_{x\to0}
\frac{1-\cos3x}{x^2}.
\]
\end{exercise}

\begin{exercise}[3]
Use power reduction to integrate

\[
\int\sin^2x\,dx.
\]
\end{exercise}

\begin{exercise}[3]
Use product-to-sum to integrate

\[
\int\sin3x\cos2x\,dx.
\]
\end{exercise}

\begin{exercise}[3]
Show that

\[
\int_{-\pi}^{\pi}
\sin(mx)\cos(nx)\,dx=0
\]

for positive integers \(m,n\).
\end{exercise}

\begin{exercise}[4]
Prove the sine and cosine addition identities geometrically.
\end{exercise}

\begin{exercise}[4]
Prove the addition identities using rotation matrices.
\end{exercise}

\begin{exercise}[4]
Prove the addition identities using Euler's formula.
\end{exercise}

\begin{exercise}[4]
Derive general multiple-angle formulas using De Moivre's theorem.
\end{exercise}

\begin{exercise}[4]
Study the relation

\[
T_n(\cos x)=\cos(nx)
\]

for Chebyshev polynomials.
\end{exercise}

\begin{exercise}[4]
Derive orthogonality relations for sine and cosine on \([-\pi,\pi]\).
\end{exercise}

\begin{exercise}[4]
Use trigonometric identities to derive Fourier coefficients for a simple
periodic function.
\end{exercise}

\begin{exercise}[4]
Use the tangent half-angle substitution to evaluate a rational trigonometric
integral.
\end{exercise}

\begin{exercise}[4]
Study the complex extensions of sine and cosine.
\end{exercise}

%%%%%%%%%%%%%%%%%%%%%%%%%%%%%%%%%%%%%%%%%%%%%%%%%%%%%%%%%%%%
\subsection{Conceptual Summary}
\label{subsec:trigonometric-identities-summary}
%%%%%%%%%%%%%%%%%%%%%%%%%%%%%%%%%%%%%%%%%%%%%%%%%%%%%%%%%%%%

The foundational identities are

\[
\boxed{
\sin^2x+\cos^2x=1,
}
\]

\[
\boxed{
1+\tan^2x=\sec^2x,
}
\]

and

\[
\boxed{
1+\cot^2x=\csc^2x.
}
\]

Reciprocal and quotient identities reduce all six trigonometric functions to
sine and cosine.

Sum and difference identities are

\[
\boxed{
\sin(A\pm B)
=
\sin A\cos B
\pm
\cos A\sin B,
}
\]

and

\[
\boxed{
\cos(A\pm B)
=
\cos A\cos B
\mp
\sin A\sin B.
}
\]

Double-angle formulas include

\[
\boxed{
\sin2x=2\sin x\cos x
}
\]

and

\[
\boxed{
\cos2x
=
\cos^2x-\sin^2x.
}
\]

Power reduction gives

\[
\boxed{
\sin^2x
=
\frac{1-\cos2x}{2},
}
\]

\[
\boxed{
\cos^2x
=
\frac{1+\cos2x}{2}.
}
\]

Product-to-sum and sum-to-product identities connect trigonometric products
with harmonic sums.

Euler's formula,

\[
\boxed{
e^{ix}
=
\cos x+i\sin x,
}
\]

provides a unified complex-exponential representation.

The most important practical rules are:

\[
\boxed{
\text{respect the domain},
}
\]

\[
\boxed{
\text{track inverse-trig principal ranges},
}
\]

and

\[
\boxed{
\text{do not treat nonidentities as algebraic shortcuts}.
}
\]

%%%%%%%%%%%%%%%%%%%%%%%%%%%%%%%%%%%%%%%%%%%%%%%%%%%%%%%%%%%%
\subsection{Section Connections}
\label{subsec:trigonometric-identities-connections}
%%%%%%%%%%%%%%%%%%%%%%%%%%%%%%%%%%%%%%%%%%%%%%%%%%%%%%%%%%%%

\chapterconnection{
Trigonometric Identities consolidates the relationships used throughout
Precalculus, Geometry, Calculus, Differential Equations, Fourier Analysis,
Complex Analysis, Mathematical Physics, and applied modeling. Pythagorean
identities connect directly with the unit circle; addition identities connect
with rotations and complex exponentials; power-reduction and product-to-sum
identities support integration and Fourier analysis; inverse-trigonometric
relations support calculus and coordinate transformations; and
amplitude-phase identities support oscillations, waves, signal models, and
control theory. The next reference section, Limit Laws, collects the standard
rules for evaluating finite, infinite, one-sided, sequence, and asymptotic
limits.
}

%%%%%%%%%%%%%%%%%%%%%%%%%%%%%%%%%%%%%%%%%%%%%%%%%%%%%%%%%%%%
% SECTION SOURCE TRACKING
%%%%%%%%%%%%%%%%%%%%%%%%%%%%%%%%%%%%%%%%%%%%%%%%%%%%%%%%%%%%

%------------------------------------------------------------
% 14.7 Trigonometric Identities
%------------------------------------------------------------
%
% Source basis:
%   - Standard trigonometric identities.
%   - Standard unit-circle, complex-exponential, inverse-trigonometric,
%     calculus, Fourier, and harmonic identities.
%
% Used for:
%   - six trigonometric functions;
%   - reciprocal identities;
%   - quotient identities;
%   - Pythagorean identities;
%   - even and odd identities;
%   - periodicity;
%   - cofunction identities;
%   - quadrant and reference-angle reductions;
%   - sum and difference identities;
%   - double-angle identities;
%   - half-angle identities;
%   - power reduction;
%   - triple-angle identities;
%   - product-to-sum identities;
%   - sum-to-product identities;
%   - auxiliary-angle forms;
%   - amplitude-phase forms;
%   - Euler's formula;
%   - De Moivre's formula;
%   - Chebyshev connections;
%   - inverse trigonometric functions;
%   - principal-value ranges;
%   - inverse-trig compositions;
%   - tangent half-angle substitution;
%   - rational parameterization of the unit circle;
%   - sinusoidal transformations;
%   - tangent transformations;
%   - small-angle approximations;
%   - trigonometric limits;
%   - derivative and integral connections;
%   - Fourier orthogonality;
%   - harmonic oscillation;
%   - common identity errors;
%   - applications and exercises.
%
% Notes:
%   - Limit Laws is developed next in Section 14.8.
%   - Derivative Tables follows in Section 14.9.
%
%%%%%%%%%%%%%%%%%%%%%%%%%%%%%%%%%%%%%%%%%%%%%%%%%%%%%%%%%%%%